\subsection{Cable-এর সংখ্যা, Connector ও দূরত্ব}
\label{ch2:wired-numeric}
\begin{multicols}{2}
\textbf{Ethernet UTP Cable-এ কয়টি তার থাকে?} প্রচলিত চার-পেয়ার UTP-তে ৪ জোড়া বা ৮টি অন্তরক আবরণযুক্ত তামার তার থাকে। জোড়াগুলোর রং সাধারণত কমলা, সবুজ, নীল ও বাদামি। প্রতিটি রঙের তারের সঙ্গে সাদা ডোরাযুক্ত সঙ্গী থাকে। পেঁচানো জোড়া বাইরের নয়েজ ও পারস্পরিক ক্রসটক কমায়।\par
\textbf{RJ45 ও RJ11 কি একই Connector?} Ethernet-এ প্রচলিত ৮টি সংযোগবিন্দুর কানেক্টরকে সাধারণত RJ45 বলা হয়। এর গঠন 8P8C। টেলিফোন লাইনে প্রচলিত ছোট কানেক্টর হলো RJ11। কানেক্টরের নাম এক হলেই সব কেবলের গতি সমান হয় না।\par
\end{multicols}
\textbf{UTP ও STP:} UTP-তে অতিরিক্ত ধাতব সুরক্ষা আবরণ নেই। STP-তে তারের জোড়া বা পুরো কেবলে এই আবরণ থাকে। সুরক্ষা আবরণ ও সঠিক গ্রাউন্ডিং EMI কমায়, তবে খরচ ও স্থাপনের জটিলতা বাড়ায়। প্রচলিত টুইস্টেড-পেয়ার Ethernet চ্যানেল প্রায় ১০০ m পর্যন্ত নকশা করা হয়। কেবলের শ্রেণি, মান, প্রান্তের সংযোগ এবং যন্ত্রের ডেটা হার মিলতে হবে। Cat 3-এর ঐতিহাসিক ব্যবহার ১০ Mbps, Cat 5-এর ১০০ Mbps এবং Cat 5e-এর Gigabit Ethernet পরিচিত উদাহরণ। সব টুইস্টেড পেয়ার কেবলের গতি এক নয়।

\textbf{১০০ m কীভাবে গোনা হয়?} প্রচলিত Structured Cabling-এ প্রায় ৯০ m স্থায়ী সংযোগ ও মোট ১০ m Patch Cord মিলিয়ে ১০০ m চ্যানেল হয়। এই দৈর্ঘ্য একটি তামার সংযোগের; পুরো ক্যাম্পাস নেটওয়ার্কের নয়। সমর্থিত Cat 6A চ্যানেলে ১০ Gbps গতির সংযোগও ১০০ m পর্যন্ত পরিকল্পনা করা যায়। এজন্য কেবল, প্রান্তের সংযোগ এবং দুই প্রান্তের পোর্ট সামঞ্জস্যপূর্ণ হতে হবে।

\subsection{Coaxial-এর Rate, Thinnet ও Thicknet}
\begin{keypointbox}{Rate-এর দুটি আলাদা ঐতিহাসিক প্রসঙ্গ}
নির্দিষ্ট কোঅ্যাক্সিয়াল ব্যবস্থার প্রচলিত উদাহরণে \textbf{২০০ Mbps} হার দেওয়া হয়। নিচের \textbf{10BASE2/10BASE5 Ethernet-এর হার ১০ Mbps}। কেবলের গঠন ও নির্দিষ্ট যোগাযোগব্যবস্থার ডেটা হার আলাদা বিষয়। তাই ২০০ Mbps-কে সব কোঅ্যাক্সিয়াল সংযোগের সর্বোচ্চ ক্ষমতা বলা যায় না।
\end{keypointbox}
\begin{tabularx}{\linewidth}{|p{.22\linewidth}|X|X|}
\hline
\textbf{বিষয়} & \textbf{Thinnet} & \textbf{Thicknet}\\\hline
Ethernet নাম & 10BASE2 & 10BASE5\\\hline
Cable & পাতলা 50-ohm Coaxial & মোটা 50-ohm Coaxial\\\hline
পুরোনো Cable-এর আনুমানিক ব্যাস & ০.২৫ inch; প্রায় ৬ mm & ০.৫ inch; প্রায় ১২ mm\\\hline
Data Rate & ১০ Mbps & ১০ Mbps\\\hline
এক Segment-এর সীমা & ১৮৫ m; নামের 2 প্রায় ২০০ m বোঝায় & ৫০০ m\\\hline
ঐতিহাসিক সংযোগ & BNC ও T-Connector & Tap/Transceiver; মোটা Backbone\\\hline
\end{tabularx}
10 মানে ১০ Mbps এবং BASE মানে Baseband Signaling। কেবল টিভিতে সাধারণত 75-ohm কোঅ্যাক্সিয়াল ও F-type কানেক্টর ব্যবহৃত হয়। এটি 10BASE2 Ethernet-এর BNC সংযোগ থেকে আলাদা।

\subsection{Fiber-এর Mode, Core ও আলোর উৎস}
\textbf{ফাইবার সংযোগের নির্দিষ্ট উদাহরণ:} ১০ Gbps Ethernet-এ 10GBASE-SR মডিউল ও OM3 Multimode Fiber দিয়ে প্রায় ৩০০ m এবং OM4 দিয়ে প্রায় ৪০০ m সংযোগের মান নির্ধারিত আছে। 10GBASE-LR ও উপযুক্ত Single-mode Fiber দিয়ে প্রায় ১০ km সংযোগের মান পাওয়া যায়। এগুলো নির্দিষ্ট মডিউল ও ফাইবারের সীমা; সব ফাইবারের সর্বোচ্চ দূরত্ব এক নয়। ডেটা হার সমান হলেও মোড ও যন্ত্র বদলালে কার্যকর দূরত্ব বদলায়।

\textbf{আলোক যোগাযোগের ইতিহাস:} Charles K. Kao ও George Hockham ১৯৬৬ সালে কম ক্ষয়যুক্ত কাচের ফাইবার দিয়ে দূরের যোগাযোগের সম্ভাবনা ব্যাখ্যা করেন। Kao-কে অপটিক্যাল ফাইবার যোগাযোগের জনক বলা হয়। এই অবদানের জন্য তিনি ২০০৯ সালে পদার্থবিজ্ঞানে নোবেল পুরস্কার পান। গবেষণা, কার্যকর কেবল তৈরি ও বাজারে ব্যবহার পৃথক সময়ে বিকশিত হয়েছে।

\textbf{আলোর উৎসের নাম:} LED = Light Emitting Diode; LASER = Light Amplification by Stimulated Emission of Radiation। আলোর উৎস, ফাইবারের মোড ও গ্রাহক যন্ত্রের সামঞ্জস্যে সংযোগ তৈরি হয়। LED বা লেজার আলোর উৎস; কেবল নয়। প্রেরক ও গ্রাহকের রূপান্তরের উদাহরণ পরের বাস্তব সমস্যায় আছে।
\textbf{আলোর তরঙ্গদৈর্ঘ্য:} পরীক্ষায় প্রচলিত আনুমানিক মান হিসেবে \textbf{প্রায় ১৫০০ nm} পাওয়া যায়। প্রচলিত নিকট-অবলোহিত আলোর ব্যবহৃত পরিসর প্রায় ৮৫০, ১৩১০ ও ১৫৫০ nm। Multimode আলোচনায় ১৩০০ nm-ও ব্যবহৃত হয়। ব্যবস্থাভেদে তরঙ্গদৈর্ঘ্য ও ডেটা হার বদলায়।
\chTwoPair{\textbf{Step-index Fiber:} কোরের মধ্যে প্রতিসরণাঙ্ক প্রায় সমান থাকে। কোর থেকে ক্ল্যাডিংয়ে গেলে প্রতিসরণাঙ্ক হঠাৎ কমে। Multimode Step-index-এ আলোকপথগুলোর দৈর্ঘ্য ভিন্ন। ফলে একই আলোক স্পন্দনের অংশগুলো পৃথক সময়ে পৌঁছে ছড়িয়ে পড়ে।}{\textbf{Graded-index Fiber:} কোরের কেন্দ্র থেকে বাইরের দিকে প্রতিসরণাঙ্ক ধীরে কমে। বাইরের দীর্ঘ পথে আলো তুলনামূলক দ্রুত চলে। এতে Multimode-এর বিভিন্ন পথে বিলম্বের পার্থক্য কমে। Single/Multimode দিয়ে মোডের সংখ্যা এবং Step/Graded দিয়ে প্রতিসরণাঙ্কের বিন্যাস বোঝানো হয়।}
\textbf{অধাতব আলোকপথ:} কাচ বা প্লাস্টিকের ফাইবার বিদ্যুৎ বহনকারী ধাতব তার নয়। আলো দিয়ে ডেটা বহন করায় বাইরের EMI ফাইবারের সংকেতে সরাসরি প্রভাব ফেলে না। তবে সুরক্ষার জন্য কোনো ফাইবার কেবলের বাইরের শক্ত আবরণে ধাতু থাকতে পারে। আলোকপথ অধাতব হওয়া মানে প্রতিটি ফাইবার কেবলের সব অংশ ধাতুমুক্ত নয়।

\textbf{সুবিধার কারণ:} উচ্চ বহনক্ষমতা, কম সংকেত ক্ষয় এবং EMI-মুক্ত আলোকপথ ফাইবারকে প্রধান সংযোগে উপযোগী করে। তবে ফাইবারে সংকেত ক্ষয় শূন্য নয়। সব তামার সংকেতের তুলনায় ফাইবারে আলো ছড়িয়ে পৌঁছানোর গতি সবসময় বেশি—এমনও নয়। তীক্ষ্ণ বাঁক, কানেক্টর ও জোড়ার কারণে ক্ষয় বিবেচনা করতে হয়। যেমন, ২ km দূরের ভবনে সংযোগ দিতে উপযুক্ত ফাইবার মোড, দুই প্রান্তের ট্রান্সসিভার ও জোড়া লাগানোর ব্যবস্থা নির্বাচন করতে হয়। পুরো পথে সংকেতের শক্তি ও ক্ষয়ের হিসাব বা Link Budget-ও মেলাতে হয়। শুধু ফাইবার ব্যবহার করলেই যেকোনো দূরত্বে যেকোনো ডেটা হার নিশ্চিত হয় না।
